Deciding whether two supplier records belong to the same corporate family is a
prerequisite for spend consolidation, credit exposure aggregation, and sanctions
screening, yet it is usually evaluated as a special case of entity matching, in which
two records refer to the same real-world entity. The two tasks are not the same. A
family link connects records that are deliberately \emph{different} entities, and the
evidence for the link frequently appears in neither record. We introduce
\textsc{CorpFam}, a public benchmark of \BenchPairs candidate pairs over \BenchFamilies
corporate families, derived from \ArchiveRows United States federal award records in
which every supplier self-reports its ultimate parent to a government registry. Every
pair is stratified by \emph{name visibility}: whether the two names are identical after
normalisation, share a distinctive token, or share none at all. Because the strata have
positive rates from \BaseVisibleRate to \BaseIdenticalRate, we report per-stratum
\emph{recall}, which is base-rate invariant, rather than $F_1$, which is not.
The result is stark: the strongest of \NumMethods matchers recovers
\BestRecallIdentical of identical pairs and \BestRecallInvisible of invisible ones, and
no method exceeds \AnyInvisibleRecallMax on the latter. More consequentially, the
failure begins before matching. Blocking decides which pairs a matcher ever sees. Across
\BlkNumSchemes schemes over the full recipient roster, among them phonetic keys,
attribute keys that ignore the name, and semantic nearest neighbours, none a function of
the token overlap that defines the stratum, no scheme exceeds three percent on invisible
pairs, and their union recovers \BlkUnionInvisible. Roughly
\BlkInvisibleLost of these links are absent from the candidate set before any matcher
runs, so no improvement at the matching stage can reach them. The links are nonetheless
real: against SEC Exhibit~21 subsidiary schedules, filed under securities law and
sharing no provenance with procurement registration, \EdgarInvisibleConfirmPct of
invisible links are corroborated, against \EdgarRandomParentPct when the same procedure
is run with permuted parents. We argue that corporate-family resolution is a retrieval
problem misfiled as a matching problem, and that the intervention point is candidate
generation rather than ranking. The benchmark, the full adjudication log, and code
reproducing every number are released.
